%% ============================================================
%% PARTE V — GUIA PARA O PROFESSOR
%% ============================================================

\chapter{Organiza\c{c}\~{a}o do Tempo de Ensino}

%% ----------------------------------------------------------------
\section{Cronograma Semanal Sugerido}
%% ----------------------------------------------------------------

O cronograma abaixo apresenta a distribuição semanal das atividades
no m\'{e}todo h\'{i}brido Fonico-Kumon com T\'{e}cnica da Boquinha.
Cada atividade possui frequ\^{e}ncia e dura\c{c}\~{a}o definidas para
otimizar o tempo de sala de aula.

\begin{table}[h]
\centering
\caption{Cronograma Semanal --- 1º Ano do Ensino Fundamental}
\label{tab:cronograma}
\resizebox{\textwidth}{!}{%
\begin{tabular}{|l|c|c|p{4cm}|}
\hline
\textbf{Atividade} & \textbf{Frequ\^{e}ncia} & \textbf{Dura\c{c}\~{a}o} & \textbf{Descri\c{c}\~{a}o} \\
\hline
Apresenta\c{c}\~{a}o de nova letra (Boquinha) & 2x por semana & 20 min &
Posi\c{c}\~{a}o da boca, tra\c{c}ado, som da letra; uso do projetor
ou quadro. \\
\hline
Pr\'{a}tica Kumon (folhas) & Di\'{a}ria & 15--20 min &
Exerc\'{i}cios de progress\~{a}o granular em folhas individuais;
cada aluno avan\c{c}a no seu ritmo. \\
\hline
Revis\~{a}o semanal & 1x por semana & 30 min &
Revis\~{a}o coletiva do conte\'{u}do da semana; jogos e atividades
l\'{u}dicas de fixa\c{c}\~{a}o. \\
\hline
Avalia\c{c}\~{a}o diagn\'{o}stica & A cada 4 semanas & 40 min &
Aplica\c{c}\~{a}o de uma das avalia\c{c}\~{o}es do Cap\'{i}tulo IV
(R-04, R-05, R-06 ou R-07). \\
\hline
Leitura em voz alta (professor) & Di\'{a}ria & 10 min &
Leitura de trechos de livros com imagens; conversa sobre a hist\'{o}ria;
exposi\c{c}\~{a}o a vocabul\'{a}rio novo. \\
\hline
Atividade l\'{u}dica (jogos) & 2x por semana & 20 min &
Jogos de mem\'{o}ria, bingo de s\'{i}labas, ca\c{c}a-palavras
e outros jogos educacionais. \\
\hline
\end{tabular}}
\end{table}

%% ----------------------------------------------------------------
\section{Sequ\^{e}ncia de Conte\'{u}do por Semana}
%% ----------------------------------------------------------------

A sequ\^{e}ncia abaixo cobre 12 semanas letivas, correspondentes ao
1º bimestre e in\'{i}cio do 2º bimestre. O professor deve ajustar
o ritmo conforme o desempenho do grupo.

\begin{table}[h]
\centering
\caption{Sequ\^{e}ncia de Conte\'{u}do --- 12 Semanas}
\label{tab:sequencia}
\begin{tabular}{|c|p{8cm}|c|}
\hline
\textbf{Semana} & \textbf{Conte\'{u}do} & \textbf{Tipo} \\
\hline
1--2 & Vogais A e E (Boquinhas, exerc\'{i}cios Kumon, jogos de rima) &
Novo \\
\hline
3--4 & Vogais I, O e U (Boquinhas, exerc\'{i}cios Kumon) &
Novo \\
\hline
5 & Consolida\c{c}\~{a}o das vogais (Revis\~{a}o semanal; avalia\c{c}\~{a}o R-04) &
Revis\~{a}o \\
\hline
6--7 & Consonantes M e S (Boquinhas, exerc\'{i}cios Kumon, palavras com M e S) &
Novo \\
\hline
8--9 & Consonantes T, R e L (Boquinhas, exerc\'{i}cios Kumon) &
Novo \\
\hline
10 & Consolida\c{c}\~{a}o do Bloco 1 (Revis\~{a}o semanal; avalia\c{c}\~{a}o R-05) &
Revis\~{a}o \\
\hline
11 & S\'{i}labas diretas com A e E; separa\c{c}\~{a}o sil\'{a}bica &
Novo \\
\hline
12 & S\'{i}labas diretas com I, O e U; avalia\c{c}\~{a}o R-06 &
Novo \\
\hline
\end{tabular}
\end{table}

%% ============================================================

\chapter{Materiais Necess\'{a}rios}

%% ----------------------------------------------------------------
\section{Material do Professor}
%% ----------------------------------------------------------------

\begin{enumerate}[leftmargin=2cm]
  \item Este livro did\'{a}tico (exemplar do professor)
  \item Caderno de planejamento semanal
  \item Projetor multim\'{i}dia (para boquinhas digitais e v\'{i}deos)
  \item Quadro verde ou branco com giz/marcador colorido
  \item Cartolinas para cartilhas de parede (alfabeto e s\'{i}labas)
  \item Tesoura e cola para recortes e montagens
  \item Impress\~{o}es das folhas de pr\'{a}tica Kumon
  \item Cartazes com o alfabeto em mai\'{u}sculas e min\'{u}sculas
  \item Flashcards de letras e s\'{i}labas (comprados ou confeccionados)
  \item Pendrive com v\'{i}deos de boquinhas e \'{a}udios de pron\'{u}ncia
\end{enumerate}

%% ----------------------------------------------------------------
\section{Material do Aluno}
%% ----------------------------------------------------------------

\begin{enumerate}[leftmargin=2cm]
  \item Este livro did\'{a}tico (exemplar do aluno)
  \item Caderno de caligrafia (pauta larga, com linha tracejada)
  \item L\'{a}pis grafite (HB) — m\'{i}nimo de 2 unidades
  \item Borracha macia
  \item L\'{a}pis de cor (caixa com 12 cores)
  \item Folhas de pr\'{a}tica Kumon (reproduzidas pelo professor)
  \item Caderno de registro de progresso (providenciado pela escola)
  \item Capa para o livro did\'{a}tico
\end{enumerate}

%% ----------------------------------------------------------------
\section{Material Complementar}
%% ----------------------------------------------------------------

\begin{enumerate}[leftmargin=2cm]
  \item Cartilha de parede com o alfabeto (mai\'{u}sculas e min\'{u}sculas)
  \item Flashcards de letras e s\'{i}labas (m\'{i}nimo de 30 unidades)
  \item Jogos de associa\c{c}\~{a}o fonema--grafema (mem\'{o}ria, bingo)
  \item \'{A}udios de pron\'{u}ncia das letras (dispon\'{i}veis online)
  \item V\'{i}deos de boquinhas (projetor ou tela)
  \item Massinha de modelar ou letras de EVA para atividades manipulativas
  \item Revistas e jornais velhos para recorte de letras e palavras
\end{enumerate}

%% ============================================================

\chapter{Estrat\'{e}gias de Diferencia\c{c}\~{a}o}

%% ----------------------------------------------------------------
\section{Para Alunos que Precisam de Mais Apoio}
%% ----------------------------------------------------------------

Alunos em n\'{i}vel pr\'{e}-sil\'{a}bico ou sil\'{a}bico inicial
requerem interven\c{c}\~{a}o diferenciada:

\begin{enumerate}[leftmargin=2cm]
  \item \textbf{Repeti\c{c}\~{a}o intensiva}: usar a t\'{e}cnica da
  boquinha com mais repeti\c{c}\~{a}o (5x em vez de 3x) para cada letra.

  \item \textbf{Folhas de refor\c{c}o}: fornecer exerc\'{i}cios
  adicionais do mesmo n\'{i}vel, focando nas habilidades que o aluno
  ainda n\~{a}o dominou.

  \item \textbf{Pequenos grupos}: trabalhar em grupos de 2--3 alunos,
  permitindo aten\c{c}\~{a}o mais individualizada.

  \item \textbf{Concreto para abstrato}: usar materiais manipulativos
  (letras de EVA, massinha de modelar, tiles de Scrabble) antes de
  avan\c{c}ar para a escrita no caderno.

  \item \textbf{Tempo estendido}: permitir mais tempo para cada atividade
  e n\~{a}o pressionar o aluno a acompanhar o ritmo do grupo.

  \item \textbf{Refus\~{a}o}: retornar a n\'{i}veis anteriores quando
  necess\'{a}rio, sem constrangimento; o erro \'{e} parte do processo.

  \item \textbf{Registro visual}: usar um painel individual com adesivos
  ou carinhas para que o aluno visualize seu pr\'{o}prio progresso.
\end{enumerate}

%% ----------------------------------------------------------------
\section{Para Alunos Avan\c{c}ados}
%% ----------------------------------------------------------------

Alunos em n\'{i}vel sil\'{a}bico-alfab\'{e}tico ou alfab\'{e}tico
avan\c{c}ado precisam de desafios adicionais:

\begin{enumerate}[leftmargin=2cm]
  \item \textbf{Adiantamento}: propor exerc\'{i}cios do n\'{i}vel
  seguinte (s\'{i}labas compostas, s\'{i}labas indiretas).

  \item \textbf{Escrita espont\^{a}nea}: incentivar a cria\c{c}\~{a}o
  de textos curtos (hist\'{o}rias, listas, cartas).

  \item \textbf{Desafios}: propor palavras mais complexas e frases
  elaboradas; incluir palavras com digrafos (CH, LH, NH).

  \item \textbf{Leitura aut\^{o}noma}: oferecer livros de n\'{i}vel
  adequado para leitura independente; criar um cantinho de leitura.

  \item \textbf{Mentoria}: convidar alunos avan\c{c}ados para ajudar
  colegas em atividades de dupla, refor\c{c}ando o aprendizado de ambos.

  \item \textbf{Projeto pessoal}: orientar o aluno a criar um mini-livro
  ilustrado com palavras e frases que ele mesmo produz.
\end{enumerate}

%% ----------------------------------------------------------------
\section{Para Alunos com Defici\^{e}ncia ou Necessidades Especiais}
%% ----------------------------------------------------------------

Alunos com defici\^{e}ncia ou necessidades educativas especiais
requerem adapta\c{c}\~{o}es curriculares e recursos assistivos:

\begin{enumerate}[leftmargin=2cm]
  \item \textbf{Adapta\c{c}\~{o}es curriculares}: adaptar exerc\'{i}cios
  conforme a defici\^{e}ncia e as orienta\c{c}\~{o}es do PAEE
  (Plano de Atendimento Educacional Especializado).

  \item \textbf{Recursos assistivos}: usar material ampliado (letras em
  alto relevo, braille), \'{a}udiogravadores, softwares de
  comunica\c{c}\~{a}o alternativa.

  \item \textbf{Apoio especializado}: incluir orienta\c{c}\~{a}o de
  profissionais da educa\c{c}\~{a}o especial (fonoaudi\'{o}logo,
  psicopedagogo, terapeuta ocupacional).

  \item \textbf{Fam\'{i}lia}: envolver os pais ou respons\'{a}veis no
  processo, com orienta\c{c}\~{a}o espec\'{i}fica para atividades em casa.

  \item \textbf{Tempo e formato}: oferecer mais tempo, provas com menos
  itens, formatos alternativos (orais, com figuras, digitais).

  \item \textbf{Avalia\c{c}\~{a}o adaptada}: utilizar as mesmas
  avalia\c{c}\~{o}es diagn\'{o}sticas, mas com adapta\c{c}\~{o}es
  previstas no Plano Educacional Individualizado (PEI).
\end{enumerate}

%% ============================================================

\chapter{Registro de Progresso}

%% ----------------------------------------------------------------
\section{Caderno Individual do Aluno}
%% ----------------------------------------------------------------

Manter um caderno individual para cada aluno \'{e} essencial para
acompanhar a evolu\c{c}\~{a}o e planejar interven\c{c}\~{o}es
pedag\'{o}gicas. O caderno deve conter:

\begin{enumerate}[leftmargin=2cm]
  \item \textbf{Dados do aluno}: nome completo, data de nascimento,
  turma, n\'{u}mero de matr\'{i}cula, nome dos respons\'{a}veis.

  \item \textbf{Data de cada atividade}: quando foi realizada a atividade
  (fólio de pr\'{a}tica, avalia\c{c}\~{a}o, jogo, etc.).

  \item \textbf{Nota obtida}: em cada folha de pr\'{a}tica, com
  indica\c{c}\~{a}o do n\'{u}mero de acertos e erros.

  \item \textbf{Observa\c{c}\~{o}es do professor}: dificuldades
  espec\'{i}ficas, avan\c{c}os observados, comportamento,
  estrat\'{e}gias utilizadas pelo aluno.

  \item \textbf{N\'{i}vel de alfabetiza\c{c}\~{a}o atual}: pré-silábico,
  silábico, silábico-alfabético ou alfabético.

  \item \textbf{Meta para o pr\'{o}ximo per\'{i}odo}: o que deve ser
  trabalhado nas pr\'{o}ximas semanas, com base na avalia\c{c}\~{a}o
  diagn\'{o}stica.
\end{enumerate}

%% ----------------------------------------------------------------
\section{Modelo de Registro}
%% ----------------------------------------------------------------

\begin{table}[h]
\centering
\caption{Modelo de Registro de Progresso Individual}
\label{tab:registro-modelo}
\resizebox{\textwidth}{!}{%
\begin{tabular}{|p{1.8cm}|p{3.5cm}|p{1.5cm}|p{1.5cm}|p{4cm}|}
\hline
\textbf{Data} & \textbf{Atividade} & \textbf{Nota} & \textbf{N\'{i}vel} & \textbf{Observa\c{c}\~{a}o} \\
\hline
01/03 & Folha 1.1 (Letra A) & 6/6 & Pr\'{e}-sil. & Reconhece todas as vogais \\
\hline
03/03 & Folha 1.2 (Letra E) & 5/6 & Pr\'{e}-sil. & Dificuldade com letra E \\
\hline
05/03 & Folha 1.3 (Letra I) & 6/6 & Pr\'{e}-sil. & Boa evolu\c{c}\~{a}o \\
\hline
08/03 & Revis\~{a}o Semanal 1 & 9/10 & Sil. & Consolida\c{c}\~{a}o adequada \\
\hline
12/03 & Avalia\c{c}\~{a}o R-04 & 28/40 & Sil. & N\'{i}vel sil\'{a}bico inicial \\
\hline
15/03 & Folha 2.1 (Letra M) & 4/6 & Sil. & Forma s\'{i}labas com ajuda \\
\hline
\end{tabular}}
\end{table}

%% ----------------------------------------------------------------
\section{Relat\'{o}rio de Progresso Mensal}
%% ----------------------------------------------------------------

A cada m\^{e}s, emitir um relat\'{o}rio para a fam\'{i}lia contendo:

\begin{enumerate}[leftmargin=2cm]
  \item Resumo das atividades realizadas no per\'{i}odo
  \item M\'{e}dia de notas nas folhas de pr\'{a}tica
  \item N\'{i}vel de alfabetiza\c{c}\~{a}o atual (com base na
  avalia\c{c}\~{a}o diagn\'{o}stica mais recente)
  \item Dificuldades identificadas e interven\c{c}\~{o}es adotadas
  \item Recomenda\c{c}\~{o}es para atividades em casa
  \item Pr\'{o}ximos passos e metas para o per\'{i}odo seguinte
  \item Assinatura do professor e do respons\'{a}vel
\end{enumerate}

%% ============================================================

\chapter{Atividades Complementares}

%% ----------------------------------------------------------------
\section{Jogos Educacionais}
%% ----------------------------------------------------------------

\subsection{Jogo da Mem\'{o}ria de S\'{i}labas}
\begin{itemize}[leftmargin=2cm]
  \item \textbf{Material}: cart\~{o}es com s\'{i}labas (MA, ME, MI, MO, MU)
  \item \textbf{Regra}: virar dois cart\~{o}es; se formarem par de s\'{i}labas
  que formam uma palavra, permanecem virados
  \item \textbf{Objetivo}: encontrar todos os pares e ler a palavra formada
  \item \textbf{Habilidade}: reconhecimento visual de s\'{i}labas e
  forma\c{c}\~{a}o de palavras
\end{itemize}

\subsection{Ca\c{c}a-Palavras}
\begin{itemize}[leftmargin=2cm]
  \item \textbf{Material}: grade com letras e lista de palavras
  \item \textbf{Regra}: encontrar e circular palavras na grade
  \item \textbf{Objetivo}: encontrar todas as palavras da lista
  \item \textbf{Habilidade}: reconhecimento visual de s\'{i}labas em
  contexto; aten\c{c}\~{a}o e concentra\c{c}\~{a}o
\end{itemize}

\subsection{Bingo de S\'{i}labas}
\begin{itemize}[leftmargin=2cm]
  \item \textbf{Material}: cart\~{o}es de bingo com s\'{i}labas
  \item \textbf{Regra}: professor sorteia s\'{i}labas; aluno marca no cart\~{a}o
  \item \textbf{Objetivo}: completar linha, coluna ou diagonal
  \item \textbf{Habilidade}: reconhecimento auditivo e visual de s\'{i}labas
\end{itemize}

\subsection{Jogo da Forca de S\'{i}labas}
\begin{itemize}[leftmargin=2cm]
  \item \textbf{Material}: quadro ou papel, l\'{a}pis
  \item \textbf{Regra}: professor pensa em uma palavra; alunos tentam
  adivinhar letras; a cada erro, desenha-se uma parte do corpo
  \item \textbf{Objetivo}: adivinhar a palavra antes de completar o desenho
  \item \textbf{Habilidade}: hip\'{o}tese ortogr\'{a}fica e decodifica\c{c}\~{a}o
\end{itemize}

%% ----------------------------------------------------------------
\section{Atividades em Sala de Aula}
%% ----------------------------------------------------------------

\subsection{Roda de Leitura}
\begin{itemize}[leftmargin=2cm]
  \item \textbf{Frequ\^{e}ncia}: di\'{a}ria (10 minutos)
  \item \textbf{Formato}: professor l\^{e} em voz alta; alunos acompanham
  com o livro ou com imagens projetadas
  \item \textbf{Material}: livros de imagens com texto simples, cartilhas
  \item \textbf{Objetivo}: exposi\c{c}\~{a}o ao modelo de leitura fluente;
  vocabul\'{a}rio; compreens\~{a}o oral
\end{itemize}

\subsection{Painel de Palavras}
\begin{itemize}[leftmargin=2cm]
  \item \textbf{Frequ\^{e}ncia}: semanal
  \item \textbf{Formato}: criar um painel com palavras novas da semana,
  recortadas de revistas ou escritas à m\~{a}o
  \item \textbf{Material}: cartolina, revistas velhas, tesoura sem ponta, cola
  \item \textbf{Objetivo}: expans\~{a}o de vocabul\'{a}rio; leitura em contexto;
  produ\c{c}\~{a}o coletiva
\end{itemize}

\subsection{Teatro de Bonecos}
\begin{itemize}[leftmargin=2cm]
  \item \textbf{Frequ\^{e}ncia}: quinzenal
  \item \textbf{Formato}: apresentar cenas com frases simples usando
  bonecos de m\~{a}o ou fantoches
  \item \textbf{Material}: bonecos de m\~{a}o, cen\'{a}rio simples, roupas
  \item \textbf{Objetivo}: produ\c{c}\~{a}o oral; leitura em contexto
  teatral; expressividade
\end{itemize}

\subsection{Canto das S\'{i}labas}
\begin{itemize}[leftmargin=2cm]
  \item \textbf{Frequ\^{e}ncia}: 2x por semana (5 minutos)
  \item \textbf{Formato}: cantar canções que repetem s\'{i}labas e palavras
  (ex.: ``Atirei o pau no gato, to -- a -- toi -- rei -- o -- pau -- no
  -- ga -- to -- TO'')
  \item \textbf{Objetivo}: consci\^{e}ncia fonol\'{o}gica; segmenta\c{c}\~{a}o
  sil\'{a}bica; prazer com a l\'{i}ngua
\end{itemize}

%% ============================================================

\chapter{Orienta\c{c}\~{o}es para Fam\'{i}lias}

%% ----------------------------------------------------------------
\section{Como Ajudar em Casa}
%% ----------------------------------------------------------------

A participa\c{c}\~{a}o da fam\'{i}lia \'{e} fator decisivo no processo
de alfabetiza\c{c}\~{a}o. Oriente os pais e respons\'{a}veis a:

\begin{enumerate}[leftmargin=2cm]
  \item \textbf{Leia para a crian\c{c}a}: todos os dias, pelo menos
  15 minutos. A leitura em voz alta develop vocabul\'{a}rio, consci\^{e}ncia
  fonol\'{o}gica e h\'{a}bito de leitura.

  \item \textbf{Cante can\c{c}\~{o}es}: can\c{c}\~{o}es infantis e
  parlendas ajudam na consci\^{e}ncia fonol\'{o}gica e na
  segmenta\c{c}\~{a}o sil\'{a}bica.

  \item \textbf{Pratique as letras}: peça que a crian\c{c}a escreva
  letras e palavras que estudou na escola; usar Letras de Scrabble
  ou cart\~{o}es pode tornar a atividade mais l\'{u}dica.

  \item \textbf{Visite bibliotecas}: leve a crian\c{c}a a livrarias
  e bibliotecas p\'{u}blicas; permita que ela escolha os livros.

  \item \textbf{Valorize o esfor\c{c}o}: elogie a dedica\c{c}\~{a}o
  e o progresso, n\~{a}o apenas o resultado final.

  \item \textbf{Comunique-se com a escola}: participe das reuni\~{o}es
  de pais e acompanhe o progresso por meio do caderno de registro.
\end{enumerate}

%% ----------------------------------------------------------------
\section{Atividades para Fam\'{i}lia}
%% ----------------------------------------------------------------

\begin{enumerate}[leftmargin=2cm]
  \item \textbf{Jogo do alfabeto}: encontrar letras em embalagens,
  revistas, placas de rua, outdoors; a crian\c{c}a coleciona as
  letras encontradas.

  \item \textbf{Cozinha educativa}: preparar receitas simples enquanto
  pratica palavras que contêm as letras estudadas (ex.: ``molho'',
  ``sal'', ``risoto'').

  \item \textbf{Passeios}: identificar letras em placas de rua, nomes
  de lojas, n\'{u}meros de ônibus; escrever o nome das coisas que veem.

  \item \textbf{Leitura compartilhada}: ler juntos e conversar sobre
  a hist\'{o}ria; perguntar ``O que aconteceu primeiro?'' e ``Como
  voc\^{e} acha que termina?''

  \item \textbf{Caixa de letras}: montar uma caixa com letras de EVA,
  Scrabble ou papel; a crian\c{c}a forma palavras e s\'{i}labas
  livremente.
\end{enumerate}

%% ----------------------------------------------------------------
\section{Sinais de Aten\c{c}\~{a}o}
%% ----------------------------------------------------------------

Procure a escola se a crian\c{c}a apresentar:
\begin{enumerate}[leftmargin=2cm]
  \item Dificuldade extrema em reconhecer letras ap\'{o}s 3 meses
  de pr\'{a}tica consistente
  \item Resist\^{e}ncia persistente em participar de atividades
  de leitura e escrita
  \item Confus\~{a}o frequente entre letras visualmente semelhantes
  (b/d, p/q, m/n) que n\~{a}o se resolve com interven\c{c}\~{a}o
  \item Dificuldade em seguir instru\c{c}\~{o}es simples e sequenciais
  \item Comportamento de frustra\c{c}\~{a}o excessivo durante atividades
  escolares (choro, recusa, agita\c{c}\~{a}o)
  \item Falta de interessa por histórias e livros em idade adequada
\end{enumerate}

Nesses casos, o professor deve acionar o Núcleo de Apoio Psicopedagógico
(NAP) ou a Equipe Multidisciplinar da escola para avaliação mais aprofundada.

%% ============================================================

\chapter{Alinhamento com as Diretrizes Curriculares Nacionais (DCNC)}

%% ----------------------------------------------------------------
\section{O que são as DCNC?}
%% ----------------------------------------------------------------

As Diretrizes Curriculares Nacionais do Ensino Fundamental
(BRASIL, Resolu\c{c}\~{a}o CNE/CEB nº 2/2012) estabelecem os
princ\'{i}pios, fundamentos e procedimentos para o ensino
fundamental nas redes p\'{u}blicas e privadas do Brasil.
Elas complementam a Base Nacional Comum Curricular (BNCC)
ao definir como o conte\'{u}do deve ser organizado e
ensinado em sala de aula.

%% ----------------------------------------------------------------
\section{Princ\'{i}pios das DCNC Aplicados à Alfabetiza\c{c}\~{a}o}
%% ----------------------------------------------------------------

Os princ\'{i}pios fundamentais das DCNC e sua aplica\c{c}\~{a}o
no m\'{e}todo h\'{i}brido Fonico-Kumon s\~{a}o:

\begin{enumerate}[leftmargin=2cm]
  \item \textbf{Integralidade}: a avalia\c{c}\~{a}o diagn\'{o}stica
  do Cap\'{i}tulo IV avalia o aluno em todas as dimens\~{o}es
  (reconhecimento, forma\c{c}\~{a}o de s\'{i}labas, leitura e
  produ\c{c}\~{a}o escrita), respeitando a integralidade do
  processo de alfabetiza\c{c}\~{a}o.

  \item \textbf{Humanismo}: o m\'{e}tdo coloca o aluno como centro
  do processo; a progress\~{a}o granular do Kumon permite que cada
  estudante avance no seu ritmo, sem competi\c{c}\~{a}o.

  \item \textbf{Contextualiza\c{c}\~{a}o}: as atividades usam
  palavras e frases do cotidiano do aluno (``mesa'', ``casa'',
  ``sapo''), conectando o conte\'{u}do escolar à vida real.

  \item \textbf{Interdisciplinaridade}: a leitura em voz alta
  e o painel de palavras integram L\'{i}ngua Portuguesa com
  Arte, Educa\c{c}\~{a}o F\'{i}sica e outras \'{a}reas.

  \item \textbf{Diferencia\c{c}\~{a}o}: as estrat\'{e}gias do
  Cap\'{i}tulo V atendem a diferentes perfis de aprendizes,
  incluindo alunos com defici\^{e}ncia e alunos avan\c{c}ados.

  \item \textbf{Avalia\c{c}\~{a}o cont\'{i}nua}: as avalia\c{c}\~{o}es
  diagn\'{o}sticas s\~{a}o formativas e cont\'{i}nuas, n\~{a}o
  somativas, conforme preconizam as DCNC.
\end{enumerate}

%% ----------------------------------------------------------------
\section{Organiza\c{c}\~{a}o do Tempo Escolar}
%% ----------------------------------------------------------------

As DCNC determinam que o tempo escolar deve ser organizado de forma
a garantir a integralidade do ensino. No m\'{e}tdo Fonico-Kumon,
isso se traduz em:

\begin{itemize}[leftmargin=2cm]
  \item \textbf{Atividade di\'{a}ria}: a pr\'{a}tica Kumon (15--20 min)
  garante a necessidade de repeti\c{c}\~{a}o e fixa\c{c}\~{a}o
  prevista nas DCNC.

  \item \textbf{Revis\~{a}o semanal}: a revis\~{a}o coletiva (30 min)
  permite ao professor identificar lacunas e reorientar o ensino.

  \item \textbf{Avalia\c{c}\~{a}o peri\'{o}dica}: as avalia\c{c}\~{o}es
  a cada 4 semanas cumprem o papel de avalia\c{c}\~{a}o diagn\'{o}stica
  e formativa exigido pelas DCNC.

  \item \textbf{Atividades l\'{u}dicas}: os jogos e brincadeiras
  educacionais (2x por semana) respeitam a faixa et\'{a}ria e a
  necessidade de brincar para aprender.
\end{itemize}

%% ----------------------------------------------------------------
\section{Gest\~{a}o de Processos de Ensino-Aprendizagem}
%% ----------------------------------------------------------------

As DCNC exigem que a gest\~{a}o dos processos de ensino-aprendizagem
seja baseada em evid\^{e}ncias e dados. O m\'{e}tdo Fonico-Kumon
atende a essa exigência por meio de:

\begin{itemize}[leftmargin=2cm]
  \item \textbf{Registro de progresso}: o caderno individual do aluno
  (Cap\'{i}tulo V) documenta a evolu\c{c}\~{a}o de cada estudante,
  permitindo an\'{a}lise longitudinal.

  \item \textbf{Classifica\c{c}\~{a}o por n\'{i}vel}: a tabela de
  classifica\c{c}\~{a}o por n\'{i}vel de alfabetiza\c{c}\~{a}o
  fornece dados objetivos para o planejamento did\'{a}tico.

  \item \textbf{Relat\'{o}rio mensal}: o relat\'{o}rio para a fam\'{i}lia
  cumpre a fun\c{c}\~{a}o de comunica\c{c}\~{a}o pedag\'{o}gica
  prevista nas DCNC.

  \item \textbf{Interven\c{c}\~{a}o diferenciada}: as estrat\'{e}gias
  de diferencia\c{c}\~{a}o s\~{a}o planejadas com base nos dados
  coletados nas avalia\c{c}\~{o}es diagn\'{o}sticas.
\end{itemize}

%% ----------------------------------------------------------------
\section{Avalia\c{c}\~{a}o conforme as DCNC}
%% ----------------------------------------------------------------

As DCNC estabelecem que a avalia\c{c}\~{a}o deve ser:
\begin{itemize}[leftmargin=2cm]
  \item \textbf{Formativa}: acompanha o processo (avalia\c{c}\~{o}es
  a cada 4 semanas);
  \item \textbf{Diagn\'{o}stica}: identifica o n\'{i}vel atual
  (classifica\c{c}\~{a}o por n\'{i}vel de alfabetiza\c{c}\~{a}o);
  \item \textbf{Processual}: registra a evolu\c{c}\~{a}o (caderno
  individual);
  \item \textbf{Participativa}: envolve o aluno, a fam\'{i}lia e a
  comunidade escolar (relat\'{o}rio mensal).
\end{itemize}

O m\'{e}tdo h\'{i}brido Fonico-Kumon est\'{a} integralmente alinhado
aos princ\'{i}pios de avalia\c{c}\~{a}o das DCNC, garantindo que
nenhum aluno seja reprovado ou rotulado negativamente, mas sim
acompanhado em seu processo de aprendizagem.

%% ============================================================
%% NOVO (R201): QUADRO DE ROTAS INDIVIDUALIZADAS + NEUROINCLUSÃO
%% ============================================================
\chapter{Rotas Individualizadas e Neuroinclusão na Prática}

\section{Quadro de Acompanhamento das Rotas (1A--1D)}

\begin{center}
\begin{tabular}{|l|l|c|c|c|}
\hline
\textbf{Aluno} & \textbf{Diagnóstico} & \textbf{Rota} & \textbf{Nível K} & \textbf{Reavaliar em} \\
\hline
\underline{\hspace{3cm}} & \underline{\hspace{1cm}}\% & \underline{\hspace{1cm}} & \underline{\hspace{1cm}} & \underline{\hspace{1.5cm}} \\
\hline
\underline{\hspace{3cm}} & \underline{\hspace{1cm}}\% & \underline{\hspace{1cm}} & \underline{\hspace{1cm}} & \underline{\hspace{1.5cm}} \\
\hline
\underline{\hspace{3cm}} & \underline{\hspace{1cm}}\% & \underline{\hspace{1cm}} & \underline{\hspace{1cm}} & \underline{\hspace{1.5cm}} \\
\hline
\underline{\hspace{3cm}} & \underline{\hspace{1cm}}\% & \underline{\hspace{1cm}} & \underline{\hspace{1cm}} & \underline{\hspace{1.5cm}} \\
\hline
\underline{\hspace{3cm}} & \underline{\hspace{1cm}}\% & \underline{\hspace{1cm}} & \underline{\hspace{1cm}} & \underline{\hspace{1.5cm}} \\
\hline
\end{tabular}
\end{center}

\rotasindividualizadas

\section{Descrição Operacional de Cada Rota}

\subsection{Rota 1A --- Recomposição Intensiva (diagnóstico abaixo de 50\%)}

\begin{itemize}
  \item \textbf{Ponto de entrada}: Unidade 0 (prontidão) + folhas Kumon A-01 a A-05.
  \item \textbf{Sessão diária}: 15 minutos, sem cronômetro; pausas livres.
  \item \textbf{Estratégia central}: multissensorialidade --- boquinha, traçado no ar,
        massinha, letras de lixa, alfabeto móvel.
  \item \textbf{Critério de saída da Rota 1A}: reconhecer as 5 vogais nas quatro
        formas gráficas e associar cada uma ao seu som, com 90\% de acerto.
\end{itemize}

\subsection{Rota 1B --- Segmentação e Leitura Funcional (50\% a 70\%)}

\begin{itemize}
  \item \textbf{Ponto de entrada}: Unidades 4 a 7 + folhas Kumon B-01 a C-03.
  \item \textbf{Estratégia central}: sílabas diretas (CV), fusão e segmentação
        oral, palavras do cotidiano com função comunicativa real.
  \item \textbf{Critério de saída da Rota 1B}: ler e escrever palavras com
        sílabas diretas e nasais simples, com 90\% de acerto.
\end{itemize}

\subsection{Rota 1C --- Consolidação e Expansão (70\% a 85\%)}

\begin{itemize}
  \item \textbf{Ponto de entrada}: Unidades 8 e 9 + folhas Kumon C-04 a D-02.
  \item \textbf{Estratégia central}: frases curtas, compreensão literal,
        produção guiada, ampliação de vocabulário temático.
  \item \textbf{Critério de saída da Rota 1C}: ler frases com compreensão e
        produzir frases simples com apoio, com 90\% de acerto.
\end{itemize}

\subsection{Rota 1D --- Aprofundamento e Autonomia (acima de 85\%)}

\begin{itemize}
  \item \textbf{Ponto de entrada}: Unidade 10 + folhas Kumon D-03 a D-06.
  \item \textbf{Estratégia central}: textos curtos, escrita autônoma, projetos
        de leitura, transferência para outras áreas do conhecimento.
  \item \textbf{Critério de consolidação da Rota 1D}: ler pequenos textos com
        fluência e compreensão e produzir frases autônomas.
\end{itemize}


\section{Orientações Neuroinclusivas para a Sala de Aula}

\subsection{Antecipação e Previsibilidade}

\begin{itemize}
  \item Apresente a rotina de seis passos no início de cada aula e mantenha-a
        visível (cartaz ou quadro). A previsibilidade reduz a ansiedade e
        libera memória de trabalho para a tarefa.
  \item Avise transições com antecedência: ``Daqui a dois minutos, vamos
        guardar o material''.
  \item Use sempre o mesmo lugar para a pista da boca, as quatro formas da
        letra e o espaço de escrita, página após página.
\end{itemize}

\subsection{Regulação Sensorial}

\begin{itemize}
  \item Respeite a opção de pausa: a caixa ``Pausa'' não é castigo nem
        privilégio; é ferramenta de autorregulação.
  \item Ofereça alternativas táteis: letras em massinha, traçado na areia,
        letras de lixa, alfabeto móvel.
  \item Permita respostas em diferentes modalidades (oral, apontando,
        traçando no ar) enquanto a escrita não estiver consolidada.
\end{itemize}

\subsection{Instrução e Comunicação}

\begin{itemize}
  \item Uma ação por instrução: ``Circule a letra A''. Depois: ``Agora,
        sublinhe as palavras''. Nunca duas ações na mesma frase.
  \item Complemente a fala com gesto e imagem: aponte, mostre, desenhe.
  \item Verifique a compreensão pedindo demonstração, não apenas
        confirmação verbal (``me mostre como você faz'').
\end{itemize}

\subsection{Velocidade Nunca é Critério}

\begin{itemize}
  \item O tempo estimado (SCT) das folhas Kumon é apenas uma referência de
        planejamento para o professor. Não o use para premiar rapidez nem
        para avaliar o aluno.
  \item O único critério de avanço é o domínio demonstrado: 90\% ou mais de
        acertos com qualidade, quantas sessões forem necessárias.
  \item Para estudantes com dislexia, o tempo adicional é direito, não
        concessão. O foco permanece na precisão e na compreensão.
\end{itemize}

\section{Checklist do Professor --- Aula Neuroinclusiva}

\begin{center}
\begin{tabular}{|l|c|}
\hline
\textbf{Verificação} & \textbf{Sim/Não} \\
\hline
Rotina de 6 passos visível para todos & $\square$ \\
\hline
Instruções com uma ação por bloco & $\square$ \\
\hline
Opções de pausa e pista disponíveis & $\square$ \\
\hline
Quatro formas da letra apresentadas & $\square$ \\
\hline
Pista articulatória multimodal oferecida & $\square$ \\
\hline
Nenhum uso de velocidade como critério & $\square$ \\
\hline
Alto contraste e espaço adequado na página & $\square$ \\
\hline
Rota individualizada registrada para cada aluno & $\square$ \\
\hline
\end{tabular}
\end{center}

\checklistdominio{Aula com os 8 itens do checklist verificados}
